\subsection{向量函数的导数}

定义 1. 设 \(\mathbf{f}\) 是定义在一个不退化为一点的区间 \(\mathbf{I} \subset \mathbf{R}\) 上的向量函数。我们说 \(\mathbf{f}\) 在一点 \(x_0 \in \mathbf{I}\) 处可导，如果

\(\lim_{x\to x_{0},x\in I,x\neq x_{0}}\frac{\mathbf{f}(x)-\mathbf{f}(x_{0})}{x-x_{0}}\) 存在（在 f 取值的那个向量空间中）；这个

极限的值称为 f 在点 \(x_{0}\) 处的一阶导数（或简称导数），记作 \(\mathbf{f}'(x_{0})\) 或 \(\mathbf{D}\mathbf{f}(x_{0})\) 。

若 f 在点 \(x_{0}\) 可导，则 f 在任何区间 \(J \subset I\) （不退化为一点且使 \(x_{0} \in J\) ）上的限制也可导，且这限制的导数等于 \(\mathbf{f}'(x_{0})\) 。反之，设 J 是含于 I 且包含 \(x_{0}\) 相对于 I 的一个邻域的区间；若 f 在 J 上的限制在点 \(x_{0}\) 处有导数，则 f 亦然。

我们把这些性质概括为一句话：导数概念是一个局部概念。

注。*1) 在运动学中，若点 \(\mathbf{f}(t)\) 是动点在空间 \(\mathbf{R}^3\) 中在时刻 \(t\) 的位置，则 \(\frac{\mathbf{f}(t) - \mathbf{f}(t_0)}{t - t_0}\) 称为时刻 \(t_0\) 与 \(t\) 之间的平均速度，而其极限 \(\mathbf{f}'(t_0)\) 是时刻 \(t_0\) 的瞬时速度（或简称速度）（当此极限存在时）。*

\begin{enumerate}
\def\labelenumi{\arabic{enumi})}
\setcounter{enumi}{1}
\tightlist
\item
  若定义在 I 上的函数 \(\mathbf{f}\) 在一点 \(x_0 \in I\) 处可导，则它在该点必相对于 I 连续。
\end{enumerate}

定义 2. 设 \(\mathbf{f}\) 是定义在区间 \(\mathrm{I} \subset \mathbf{R}\) 上的向量函数，并设 \(x_0\) 是 \(\mathrm{I}\) 中使区间 \(\mathrm{I} \cap [x_0, +\infty[\) （相应地，\(\mathrm{I} \cap ] - \infty, x_0]\) ）不退化为一点的点。若 \(\mathbf{f}\) 在点 \(x_0\) 右可导（相应地，左可导），指的是：\(\mathbf{f}\) 在区间 \(\mathrm{I} \cap [x_0, +\infty[\) （相应地，\(\mathrm{I} \cap ] - \infty, x_0]\) ）上的限制在点 \(x_0\) 可导；这一限制在点 \(x_0\) 的导数值称为 \(\mathbf{f}\) 在点 \(x_0\) 的右导数（相应地，左导数），记作 \(\mathbf{f}_d'(x_0)\) （相应地，\(\mathbf{f}_g'(x_0)\) ）。

设 f 是定义在 I 上的向量函数，\(x_{0}\) 是 I 的一个内点且 f 在该点连续；由定义 1 和 2 可知，f 在 \(x_{0}\) 可导当且仅当 f 在该点既有右导数又有左导数，且这两个导数相等；此时

\[
\mathbf {f} ^ {\prime} (x _ {0}) = \mathbf {f} _ {d} ^ {\prime} (x _ {0}) = \mathbf {f} _ {g} ^ {\prime} (x _ {0}).
\]

例。1) 常函数在每一点的导数都为零。

\begin{enumerate}
\def\labelenumi{\arabic{enumi})}
\setcounter{enumi}{1}
\item
  仿射线性函数 \(x \mapsto ax + b\) 在每一点的导数都等于 a。
\item
  实函数 \(1 / x\) （对 \(x \neq 0\) 有定义）在每个点 \(x_0 \neq 0\) 处可导，因为我们有 \(\left(\frac{1}{x} - \frac{1}{x_0}\right) / (x - x_0) = -\frac{1}{xx_0}\) ，而且由于 \(1 / x\) 在 \(x_0\) 连续，前述表达式的极限为 \(-1 / x_0^2\) 。
\item
  数量函数 \(|x|\) 定义在 \(\mathbf{R}\) 上，它具有右导数 \(+1\) 与左导数 \(-1\) （在 \(x = 0\) 处）；它在该点不可导。
\end{enumerate}

*5) 实函数在 \(x = 0\) 时等于 0、等于 \(x \sin 1/x\) （当 \(x \neq 0\) ），它在 \(\mathbf{R}\) 上有定义且连续，但在点 \(x \neq 0\) 处既无右导数亦无左导数。可以给出在区间上连续而在区间每一点都没有导数的函数的例子（I, p.~35, exerc. 2 和 3）。

定义 3. 设 \(\mathbf{f}\) 是定义在区间 \(\mathrm{I} \subset \mathbf{R}\) 上的向量函数；若它在 \(\mathrm{I}\) 的每一点可导（相应地，右可导、左可导），则称它在 \(\mathrm{I}\) 上可导（相应地，右可导、左可导）；函数 \(x \mapsto \mathbf{f}'(x)\) （相应地，\(x \mapsto \mathbf{f}_d'(x)\) 、\(x \mapsto \mathbf{f}_g'(x)\) ）定义在 \(\mathrm{I}\) 上，称为 \(\mathbf{f}\) 的导函数，或（滥用语言的说法）导数（相应地，右导数、左导数），记作 \(\mathbf{f}'\) 或 \(\mathrm{Df}\) 或 \(df/dx\) （相应地，\(\mathbf{f}_d', \mathbf{f}_g'\) ）。

注。一个函数可以在某个区间上可导，而其导数不必在区间的每一点连续（cf.~I, p.~36, exerc. 5）；*在 x = 0 时等于 0、等于 \(x^{2} \sin 1/x\) （当 \(x \neq 0\) ）的函数的例子就表明了这一点；它处处有导数，但这导数在点 \(x = 0\) 不连续*

